\section{SHAKE-256 Binding}
\label{sec:shake256}

\tesserahash{} invokes SHAKE-256 \citep{fips202} unchanged from FIPS 202: \code{SHAKE256(content, 64 bytes)} produces 512 bits of pseudorandom output deterministic in the input. There is no salt, no domain separator, no length-prefix wrapper, and no compression layer above SHAKE-256 itself. Two callers hashing identical content produce identical 64-byte tails.

\subsection{Why SHAKE-256 over SHA3-512, SHA-512, BLAKE3, or Blake2b}

\begin{itemize}
  \item \textbf{Vs SHA3-512.} Both are Keccak-f[1600] sponge constructions with a 1024-bit output and 256-bit collision resistance. SHA3-512 \citep{fips202} pads with the SHA-3 domain separator \code{01}; SHAKE-256 pads with \code{1111}. We chose SHAKE-256 because (a) the SHAKE family's variable-output API is the one we want for any future TesseraHash variant that adjusts tail size, and (b) the rest of the \tessera{} monorepo uses SHAKE-256 (in \tesseraseal{}'s HKDF salt construction and hybrid combiner) and we already vendor and audit it once.
  \item \textbf{Vs SHA-512.} SHA-512 is a Merkle--Damg{\aa}rd construction susceptible to length-extension attacks. \tesserahash{} is not used as a MAC, so length extension is not directly exploitable, but Keccak-based primitives are the consensus PQ-era choice and we have no reason to swim against the standardization current \citep{fips202}.
  \item \textbf{Vs BLAKE3 or Blake2b.} BLAKE3 and Blake2b are faster on x86 and ARM than SHAKE-256. We accept the speed cost (\cref{sec:conformance:bench}) in exchange for: (i) NIST-standardized primitive, (ii) zero new vendored crypto code (Keccak is already shared with \tesseraseal{}), (iii) FIPS 202 alignment for FIPS-relevant deployments. TesseraHash is not on the latency-critical path of most use cases; throughput numbers in \cref{sec:conformance:bench} are sufficient.
\end{itemize}

\subsection{Implementation sources}

\begin{center}
\small
\begin{tabular}{@{}lll@{}}
\toprule
\textbf{Language} & \textbf{SHAKE-256 source}                                    & \textbf{Constant-time / audit} \\
\midrule
Python      & stdlib \texttt{hashlib.shake\_256}                              & CPython hashlib backed by OpenSSL \\
C           & vendored Keccak under \texttt{common/c/keccak.\{h,c\}}          & shared with \tesseraseal{}; reviewed in tessera-seal v1.0 \\
C++         & wraps the C Keccak source                                       & inherits C audit \\
Go          & stdlib \texttt{golang.org/x/crypto/sha3} (\code{sha3.NewShake256}) & Go stdlib reviewed \\
Rust        & RustCrypto \texttt{sha3} crate (\code{Shake256})                 & RustCrypto reviewed \\
TypeScript  & \texttt{@noble/hashes} \texttt{shake\_256}                       & const-time documented \\
Java        & vendored Keccak under \texttt{common/java/com/tessera/common/} & shared with \tesseraseal{}; reviewed \\
CPython ext & wraps the vendored C Keccak                                     & inherits C audit \\
\bottomrule
\end{tabular}
\end{center}

The cross-project sharing of Keccak under \texttt{common/} is a deliberate scope decision: \tesserahash{} alone could vendor its own copy, but \tesseraseal{} requires HMAC-SHA3-256 and HKDF-SHA3-256 for its KDF schedule, so the second consumer makes the \texttt{common/} extraction load-bearing.

\subsection{What about salts?}

Salts would defeat content-equivalence. A salted TesseraHash would no longer be a content identifier; it would be a randomized commitment. If your use case is ``commit to content without revealing it,'' you want a keyed MAC (HMAC-SHAKE-256, or KMAC256), not \tesserahash{}. \tesserahash{} is the unsalted content hash by design.

\subsection{Variable-length input}

SHAKE-256 takes arbitrary-length input and \tesserahash{} passes it through unchanged. Content larger than memory should be streamed into the Keccak sponge using each language's incremental API (\texttt{shake.update(chunk)}); the cumulative output is identical to a single-shot hash of the concatenation. The Python reference \texttt{hash\_only(content)} accepts \texttt{bytes}, \texttt{bytearray}, or \texttt{memoryview}; the C reference takes a \texttt{const uint8\_t *} content pointer and a \texttt{size\_t} length and zero-copies into the sponge.
